\section{Introduction}
\label{sec:introduction}

Sparse linear codes appear inside several pseudorandom correlation generators
(PCGs) and silent cryptographic protocols.  These applications do not begin with a
received word and ask a decoder to correct errors.  They instead need a
sampled linear map to send every nonzero message to an output with many
nonzero coordinates.  In coding terminology, the generator must have large
minimum distance.

Expand--Accumulate (EA) and Expand--Convolute (EC) are two such families
introduced for cryptographic applications~\cite{EA,EC}.  Both first apply a
sparse random matrix.  EA follows that matrix by a running sum.  EC follows it
by a randomized recursive convolution.  The sparse first stage keeps encoding
cheap, while the recursive second stage spreads each nonzero contribution
through the output.

The published distance proofs do not analyze the recursive output law
exactly.  They compare or bound a state process and then apply concentration.
This choice obscures two different regimes.  A message with linearly many
nonzero symbols makes the expander output nearly uniform; its distance is
governed by the same entropy calculation as a random linear code.  A message
with only a few nonzero symbols instead tests whether the sparse first stage
creates enough opportunities for the recursive map to become active.  The
earlier bounds lose substantially in both regimes.

This paper keeps the exact input--output weight enumerator of the recursive
map.  An enumerator is a generating function that records how many inputs of
each weight produce outputs of each weight.  Transfer matrices compute this
function while retaining only the information about the current state that
can affect future output weight.  The accumulator needs two states.  A
convolution of memory $m$ needs only $m+1$ states, rather than all $p^m$
possible field-valued states.

\paragraph{Two cases.}
We first study binary codes.  This case models bit-valued choices, as in
Silent OT, and the Hamming support determines every message value.  We then
study codes over a finite field $\F_p$ whose order is close to $2^{128}$.
This case models field-valued correlations, as in Silent VOLE.  Here equal
supports can cancel differently, so the construction and the proof must mix
field values.  The two cases share the enumerator method but have different
sparse-message bottlenecks.

\paragraph{Binary analysis of the original ensembles.}
Our first results change only the analysis.  For the binary Bernoulli EA
ensemble, write the entry probability as $C\ln(n)/n$, so an expander row has
expected weight $C\ln n$.  The exact transfer bound lowers the required
constant at relative distance $0.05$---meaning output weight $0.05n$---from
$C>4.938$ to $C>1.773$.  At rate $1/5$, the
dense-message condition improves from a distance below $0.1127$ to the binary
Gilbert--Varshamov (GV) value $0.2430$.  For wrapped EC with memory $21$, the
same method lowers the density requirement at distance $0.05$ from
$C>10.4344$ to $C>1.111623$.  The permitted rate rises from below $0.1278$ to
the GV threshold $0.7136$.

The GV value is a benchmark defined by the volume of a Hamming ball.  It is
not a finite-length upper bound.  Saying that a certificate reaches the
floored GV cutoff means that its integer distance threshold is the floor of
the corresponding asymptotic value at that length.

\paragraph{Binary regular-expander results.}
The exact analysis identifies low-support messages as the remaining cost.
We therefore replace independent edge incidence by a region-stratified
left-regular expander.  Each input coordinate has exactly one edge in each
output region.  This change is part of the construction, not merely a proof
improvement.  At rate one half, the regular binary EA certificate uses degree
$64$ instead of Bernoulli expected degree $75$.  The regular binary EC
certificate uses degree $28$ instead of Bernoulli expected degree $42.2127$,
with convolution memory $9$.  Fixing both expander degrees improves EC much
further.  A binary two-sided regular EC ensemble uses left/right degrees
$10/5$ and convolution memory $15$.  Its cutoff reaches $99.993\%$ of the
rate-one-half GV distance, with sampling failure below $2^{-32.4990}$.
Two intermediate profiles use degrees $14/7$ with memory $6$ and degrees
$18/9$ with memory $4$.  Their sampling failures are below $2^{-26.0875}$
and $2^{-21.6746}$, respectively.
A lower-degree profile uses degrees $6/3$ and memory $79$ at the same cutoff.
Its sampling failure is below $2^{-21.4222}$.

Scalar extension does preserve the minimum distance of a binary generator,
but it does not add field mixing.  In a basis of $\F_{2^{128}}$ over $\F_2$,
the resulting map is 128 independent copies of the binary map.  This is the
right structure for binary-choice Silent OT, but not the field-valued mixing
needed by Silent VOLE over $\F_{2^{128}}$.

\paragraph{Fields near 128 bits.}
Let $p$ be a prime power.  Over the finite field $\F_p$, cancellations depend
on message values as well as their support.
We attach an independent nonzero field label to each expander edge.  The label
symmetry makes the distribution of a fixed message depend only on its support.
Because nonzero scalar multiples have the same zero coordinates after every
linear map, we count each equivalence class of scalar-multiple messages once.
These classes are the one-dimensional message subspaces.
Constraint traces then group the resulting projective union before it becomes
too large.

For finite-field EC, we also sample each feedback vector from the full field.
This distribution preserves an exact reduced state based only on the number
of trailing zero outputs.  At rate one half over $\F_{2^{127}-1}$, we certify
the floored $p$-ary GV cutoff at length $2{,}097{,}150$.  The expander has
left degree $22$ and right degree $11$, and the convolution has memory $12$.
The probability of sampling a generator below the claimed distance is less
than $2^{-27.2951}$.

For the field-valued Silent VOLE setting, we separately instantiate the same
construction over $\F_{2^{128}}$.  Degrees $24/12$ with memory $5$ reach
$98.415\%$ of the rate-one-half GV distance.  The certified sampling-failure
probability is below $2^{-45.03}$.  Degrees $26/13$ with memory $4$ reach the
same distance with sampling failure below $2^{-135.20}$.  Both proved
ensembles use nonzero field labels on expander edges and full-field
coefficients in the convolution.

Keeping degrees $24/12$ and memory $5$ fixed gives a local field-size
profile.  At the corresponding field-relative cutoffs, the certified failure
exponents are $45.84$, $45.03$, $44.24$, and $38.78$ for field orders
$2^{127}$, $2^{128}$, $2^{129}$, and $2^{136}$.  The decline comes from the
increasing distance target rather than weaker field mixing at a fixed target.

\begin{table}[t]
\centering
\small
\setlength{\tabcolsep}{4pt}
\caption{Representative finite certificates.  The probability is over the
random choices used to sample the generator.  The final column reports the
distance, failure bound, or comparison relevant to each row.}
\label{tab:intro-results}
\begin{tabular}{@{}lcccll@{}}
\toprule
ensemble & field & rate & expander & recursive map & conclusion \\
\midrule
\multicolumn{6}{@{}l}{\emph{Binary case}} \\
EA & $\F_2$ & $1/5$ & Bernoulli, mean $50$ & accumulator
  & $U<2^{-20.46}$ at $\delta=0.05$ \\
wrapped EC & $\F_2$ & $1/5$ & Bernoulli, mean $46.42$ & memory $5$
  & $U<2^{-32.08}$; $10.01$-bit gain \\
regular EA & $\F_2$ & $1/2$ & left degree $64$ & accumulator
  & $99.988\%$ of GV; $14.67\%$ fewer edges \\
regular EC & $\F_2$ & $1/2$ & left degree $28$ & memory $9$
  & $99.994\%$ of GV; $33.67\%$ fewer edges \\
two-sided EC & $\F_2$ & $1/2$ & degrees $18/9$ & memory $4$
  & $99.993\%$ of GV; $U<2^{-21.674}$ \\
two-sided EC & $\F_2$ & $1/2$ & degrees $14/7$ & memory $6$
  & $99.993\%$ of GV; $U<2^{-26.087}$ \\
two-sided EC & $\F_2$ & $1/2$ & degrees $10/5$ & memory $15$
  & $99.993\%$ of GV; $U<2^{-32.499}$ \\
two-sided EC & $\F_2$ & $1/2$ & degrees $6/3$ & memory $79$
  & $99.993\%$ of GV; $U<2^{-21.422}$ \\
\addlinespace
\multicolumn{6}{@{}l}{\emph{Fields near 128 bits}} \\
two-sided EC & $\F_{2^{127}-1}$ & $1/2$ & labeled degrees $22/11$ & memory $12$
  & floored $p$-ary GV; $U<2^{-27.295}$ \\
two-sided EC & $\F_{2^{128}}$ & $1/2$ & labeled degrees $24/12$ & memory $5$
  & $98.415\%$ of GV; $U<2^{-45.03}$ \\
two-sided EC & $\F_{2^{128}}$ & $1/2$ & labeled degrees $26/13$ & memory $4$
  & $98.415\%$ of GV; $U<2^{-135.20}$ \\
\bottomrule
\end{tabular}
\end{table}

\paragraph{Contributions.}
The paper makes five contributions.

\begin{enumerate}
  \item We give a common exact-enumerator framework for sparse random maps
  followed by recursive rate-one maps.
  \item We sharpen the distance proofs for the original binary EA and wrapped
  EC ensembles, including direct comparisons with the published bounds.
  \item We analyze left-regular and two-sided regular expanders.  Binary EC
  gives four near-GV degree--memory points.
  \item We extend the construction and analysis to fields near 128 bits using
  edge labels, projective counting, and constraint traces.  A certified local
  profile separates field-size effects from changes in the distance target.
  \item We give reproducible finite certificates checked with outward-rounded
  ball arithmetic.  The certificates cover every nonzero message support.
\end{enumerate}

\paragraph{Scope.}
Our probability bounds concern the sampling of a generator.  They do not
analyze a decoder, a noise distribution, or a complete cryptographic
protocol.  Each large-field headline result fixes one field and one finite
length.  Increasing the field order also changes the field-relative GV
cutoff, so the displayed margins are not uniform in the field order.  The
finite claims are separate from our asymptotic theorem for each fixed field.

\paragraph{Related work.}
Exact input--output weight enumerators and decompositions by message weight
also appear in analyses of turbo-like codes~\cite{BazziMahdianSpielman2003}.
Block--Accumulate codes apply this method to a block-diagonal outer map
followed by recursive stages~\cite{BlockAccumulate2026}.  That work computes
finite expected weight spectra and uses them to select concrete parameters.
Its main code family changes the outer map.  We instead retain the original
EA and EC ensembles, replace their recursive-stage relaxations by exact
transfers, and then measure the effect of regular expander incidence.

Recent large-field analyses of repeat--accumulate--accumulate (RAA) codes
address a second source of slack: the union over many field assignments with
the same set of nonzero message coordinates.  Akhiani and Zhang group these
assignments before the union bound and prove constant relative distance over
sufficiently large fields~\cite{AkhianiZhang2026}.  Khabbazian sharpens the
analysis for randomly scaled RAA codes.  The proof fixes the support geometry,
identifies scalar multiples, and charges only the cancellation equations that
exceed the remaining message degrees of freedom~\cite{Khabbazian2026}.

Our support grouping in \eqref{eq:support-grouping} uses the same accounting
principle.  The EA and EC ensembles require different structural counts.
Expander collisions determine the rank of the accumulator constraints.  For
EC, the convolution state also determines which constraints are independent.
Our transfer matrices therefore count output weight and independent equations
together.  The large-field RAA results study an asymptotic regime in which the
field grows with the block length.  Our near-128-bit headline results instead
fix the field and length, then certify the complete finite union bound.

\paragraph{Organization.}
Section~\ref{sec:constructions} defines the sampled maps and the distance
event.  Section~\ref{sec:enumerator-method} develops the common enumerator
method.  Sections~\ref{sec:binary-ea} and~\ref{sec:binary-ec} treat the
original binary ensembles.  Section~\ref{sec:regular} introduces regular
expanders and gives the binary rate-one-half frontier.
Section~\ref{sec:prime-fields} gives the finite-field extension and separates
fixed-field asymptotics from cryptographic-size finite bounds.
Section~\ref{sec:certificates} describes the certified evaluation, and
Section~\ref{sec:discussion} records limitations and open questions.
